% ===== PRÉAMBULE CANONIQUE — à coller VERBATIM en tête de CHAQUE .tex =====
\documentclass[11pt,a4paper]{article}
\usepackage[utf8]{inputenc}\usepackage[T1]{fontenc}\usepackage{lmodern}
\usepackage[french]{babel}
\usepackage{amsmath,amssymb,mathtools}\usepackage{icomma}
\usepackage{siunitx}\sisetup{locale=FR}  % \qty{}{}, \si{}, \ang{}
\usepackage{xcolor}\usepackage{colortbl}
\usepackage{tikz}\usepackage{pgfplots}\pgfplotsset{compat=1.18}
\usepackage[siunitx,europeanresistors]{circuitikz} % circuits (élec.)
\usepackage[most]{tcolorbox}\usepackage{enumitem}\usepackage{array,booktabs}
\usepackage[a4paper,margin=2.2cm]{geometry}
\usetikzlibrary{arrows.meta,calc,angles,quotes,positioning,patterns,%
  backgrounds,shapes.geometric,decorations.pathmorphing,decorations.markings,intersections}
\definecolor{primary}{RGB}{222,49,99}\definecolor{primarydark}{RGB}{150,22,55}
\definecolor{defcol}{RGB}{0,114,178}\definecolor{methcol}{RGB}{52,92,148}
\definecolor{excol}{RGB}{0,135,90}\definecolor{attncol}{RGB}{200,40,55}
\tcbset{boxbase/.style={enhanced,breakable,boxrule=0.9pt,arc=2.5pt,
  left=3mm,right=3mm,top=2.3mm,bottom=2.3mm,fonttitle=\bfseries,coltitle=white}}
% --- boîtes TUTORIEL (noms EXACTS, ne pas renommer) ---
\newtcolorbox{cle}[1][]{boxbase,colback=defcol!6,colframe=defcol,
  title={\textsc{À retenir}\ifx\relax#1\relax\relax\else\textmd{ : \itshape #1}\fi}}
\newtcolorbox{methode}[1][]{boxbase,colback=methcol!7,colframe=methcol,sharp corners,
  title={\textsc{Méthode}\ifx\relax#1\relax\relax\else\textmd{ --- \itshape #1}\fi}}
\newtcolorbox{exemplebox}[1][]{boxbase,colback=excol!5,colframe=excol,
  title={\textsc{Exemple}\ifx\relax#1\relax\relax\else\textmd{ : \itshape #1}\fi}}
\newtcolorbox{piege}[1][]{boxbase,colback=attncol!6,colframe=attncol,
  title={\textsc{Piège}\ifx\relax#1\relax\relax\else\textmd{ : \itshape #1}\fi}}
\newtcolorbox{remarque}[1][]{boxbase,colback=black!4,colframe=black!45,
  title={\textsc{Remarque}\ifx\relax#1\relax\relax\else\textmd{ : \itshape #1}\fi}}
% --- boîtes EXERCICE / CORRIGÉ (noms EXACTS, auto-comptées) ---
\newtcolorbox[auto counter]{exo}[1][]{boxbase,colback=primary!4,colframe=primary,
  title={\textsc{Exercice}~\thetcbcounter\ifx\relax#1\relax\relax\else\textmd{ --- \itshape #1}\fi}}
\newtcolorbox[auto counter]{corrige}[1][]{boxbase,colback=excol!4,colframe=excol,
  title={\textsc{Corrigé}~\thetcbcounter\ifx\relax#1\relax\relax\else\textmd{ --- \itshape #1}\fi}}
\newcommand{\vv}[1]{\vec{#1}}\newcommand{\ds}{\displaystyle}
\newcommand{\R}{\mathbb{R}}\newcommand{\N}{\mathbb{N}}\newcommand{\Z}{\mathbb{Z}}
% (un \usepackage{fancyhdr}+en-tête est possible ; il est ignoré côté web)
% =========================================================================
\begin{document}
\begin{corrige}[moment algébrique de plusieurs forces]
\begin{enumerate}[label=\alph*)]
  \item Horaire compté $+$ : $\mathcal{M}_1 = +F_1 b_1 = +10\times 0{,}4 = +\qty{4}{N.m}$ ;
        $\mathcal{M}_2 = -F_2 b_2 = -20\times 0{,}15 = -\qty{3}{N.m}$.
  \item $\mathcal{M} = \mathcal{M}_1+\mathcal{M}_2 = +4-3 = +\qty{1}{N.m}$. Le résultat est positif : la
        barre tourne dans le sens \textbf{horaire}.
\end{enumerate}
\end{corrige}

\begin{corrige}[bras de levier ou distance au point d'application ?]
\begin{enumerate}[label=\alph*)]
  \item On utilise le \textbf{bras de levier} $b$ (distance perpendiculaire à la ligne d'action), pas $L$ :
        $\mathcal{M} = F\times b = 50\times 0{,}3 = \qty{15}{N.m}$.
  \item Si la force était perpendiculaire, le bras vaudrait $L$ : $\mathcal{M}=50\times 0{,}4=\qty{20}{N.m}$.
        C'est \emph{plus grand} : une force \textbf{perpendiculaire} au levier donne le moment maximal —
        c'est l'orientation la plus efficace.
\end{enumerate}
\end{corrige}

\begin{corrige}[est-ce un couple ?]
\begin{enumerate}[label=\alph*)]
  \item \textbf{Oui} : deux forces de \emph{même intensité} (\qty{6}{N}), \emph{parallèles} (verticales),
        de \emph{sens opposé}, à lignes d'action \emph{distinctes} ($x=0$ et $x=\qty{0.5}{m}$). Leur
        résultante est nulle : $6-6=\qty{0}{N}$.
  \item $\mathcal{M} = F\times d = 6\times 0{,}5 = \qty{3}{N.m}$.
\end{enumerate}
\end{corrige}

\begin{corrige}[le volant]
\begin{enumerate}[label=\alph*)]
  \item Les deux forces s'appliquent aux extrémités d'un diamètre : $d = 2r = 2\times 0{,}2 = \qty{0.4}{m}$.
  \item $\mathcal{M} = F\times d = 15\times 0{,}4 = \qty{6}{N.m}$.
\end{enumerate}
\end{corrige}

\begin{corrige}[travail et moment : même nombre, sens différent]
\begin{enumerate}[label=\alph*)]
  \item $W = F\times d = 25\times 2 = \qty{50}{J}$ (déplacement parallèle à la force).
  \item $\mathcal{M} = F\times b = 25\times 2 = \qty{50}{N.m}$ (bras perpendiculaire).
  \item \textbf{Non.} Même valeur numérique, mais grandeurs différentes : le travail est une
        \emph{énergie} (en joules), le moment mesure un \emph{effet de rotation} (en \si{\newton\metre},
        jamais en joules).
\end{enumerate}
\end{corrige}

\end{document}
